% !TEX program = xelatex
\documentclass[12pt,a4paper]{report}

\usepackage{fontspec}
\usepackage{polyglossia}
\setdefaultlanguage{russian}
\setmainfont{Times New Roman}[
  BoldFont = Times New Roman Bold,
  ItalicFont = Times New Roman Italic,
]
\setsansfont{Arial}
\setmonofont{Courier New}
\newfontfamily\cyrillicfont{Times New Roman}
\newfontfamily\cyrillicfontsf{Arial}
\newfontfamily\cyrillicfonttt{Courier New}

\usepackage{geometry}
\geometry{left=2.5cm,right=2cm,top=2.5cm,bottom=2.5cm}

\usepackage{booktabs}
\usepackage{longtable}
\usepackage{array}
\usepackage{tabularx}
\usepackage{graphicx}
\usepackage{xcolor}
\usepackage{listings}
\usepackage{hyperref}
\usepackage{float}
\usepackage{amsmath}
\usepackage{enumitem}
\usepackage{tikz}
\usetikzlibrary{shapes.geometric, arrows.meta, positioning, fit, backgrounds}

\hypersetup{
    colorlinks=true,
    linkcolor=blue!60!black,
    citecolor=green!50!black,
    urlcolor=blue!70!black,
    pdftitle={Guardrail Service --- Технический отчёт},
}

\definecolor{codebg}{RGB}{245,245,245}
\definecolor{keyword}{RGB}{0,80,160}

\lstdefinestyle{python}{
    language=Python,
    basicstyle=\ttfamily\small,
    backgroundcolor=\color{codebg},
    keywordstyle=\color{keyword}\bfseries,
    commentstyle=\color{gray},
    stringstyle=\color{teal},
    breaklines=true,
    frame=single,
    rulecolor=\color{gray!30},
    tabsize=4,
    showstringspaces=false,
}

\lstdefinestyle{json}{
    basicstyle=\ttfamily\small,
    backgroundcolor=\color{codebg},
    breaklines=true,
    frame=single,
    rulecolor=\color{gray!30},
}

\lstdefinestyle{bash}{
    language=Bash,
    basicstyle=\ttfamily\small,
    backgroundcolor=\color{codebg},
    breaklines=true,
    frame=single,
}

\title{%
  \textbf{Guardrail Service for Conversational AI}\\[0.4em]
  \large Технический отчёт по реализации сервиса модерации\\
  \large Тестовое задание: Backend / ML Engineering
}
\author{Middle/Senior Backend Engineer}
\date{\today}

\begin{document}

%----------------------------------------------------------------------
% Титульный лист
%----------------------------------------------------------------------
\begin{titlepage}
    \centering
    \vspace*{1.5cm}
    {\LARGE Москва, \the\year\par}
    \vspace{2cm}
    {\Huge\bfseries Guardrail Service for Conversational AI\par}
    \vspace{1.2cm}
    {\Large Технический отчёт\par}
    \vspace{0.8cm}
    {\large Сервис модерации ответов conversational AI\par}
    \vspace{0.5cm}
    {\normalsize Тестовое задание: проектирование и реализация\\
    production-style MVP на Python / FastAPI\par}
    \vfill
    \begin{tabular}{rl}
        \textbf{Автор:} & Backend Engineer \\
        \textbf{Стек:} & Python, FastAPI, Transformers, ChromaDB, OpenRouter \\
        \textbf{Дата:} & \today \\
    \end{tabular}
    \vspace{2cm}
\end{titlepage}

\pagenumbering{roman}
\tableofcontents
\clearpage
\pagenumbering{arabic}

%======================================================================
\chapter{Постановка задачи}
%======================================================================

\section{Описание проблемы}

Корпоративные и потребительские системы на базе больших языковых моделей (LLM)
генерируют ответы в режиме реального времени. Даже при наличии системных промптов
модель может выдавать токсичный контент, персональные данные, инструкции по причинению вреда
или поддаваться атакам jailbreak / prompt injection. После генерации ответа требуется
независимый контур проверки --- \emph{guardrail} --- который возвращает структурированное
решение для шлюза доставки сообщения пользователю.

\textbf{Guardrail Service} --- HTTP-сервис, принимающий текст ответа ассистента
(и опционально контекст диалога) и возвращающий решение модерации в машиночитаемом JSON.

\section{Зачем нужны guardrails для LLM}

\begin{itemize}[noitemsep]
    \item \textbf{Безопасность пользователей} --- блокировка self-harm, насилия, мошенничества.
    \item \textbf{Комплаенс} --- снижение регуляторных рисков (в т.\,ч.\ 152-ФЗ при обработке ПДн).
    \item \textbf{Репутационные риски} --- предотвращение утечки PII и токсичных формулировок.
    \item \textbf{Устойчивость к атакам} --- детекция попыток обхода политик (DAN, ignore instructions).
    \item \textbf{Аудит} --- объяснимые решения с категорией и уровнем риска.
\end{itemize}

\section{Категории рисков}

Сервис поддерживает 18 категорий нарушений политики и категорию \texttt{none}:

\begin{longtable}{@{}p{4.2cm}p{9.5cm}@{}}
\toprule
\textbf{Категория} & \textbf{Смысл} \\
\midrule
\endhead
hate\_speech & Разжигание ненависти, дискриминация \\
extremism & Экстремизм, радикализация \\
violence & Насилие, угрозы \\
self\_harm & Суицид, самоповреждение \\
medical\_advice & Неквалифицированные мед.\ рекомендации \\
legal\_advice & Юридические консультации вместо юриста \\
pii\_leakage & Утечка персональных / секретных данных \\
sexual\_content & Неприемлемый сексуальный контент \\
harassment & Травля, домогательства \\
fraud & Мошенничество, фишинг \\
malware & Вредоносное ПО, эксплойты \\
jailbreak\_attempts & Обход ограничений модели \\
prompt\_injection & Инъекции в промпт \\
political\_manipulation & Манипуляции выборами \\
disinformation & Дезинформация \\
toxic\_language & Токсичная лексика \\
privacy\_violations & Нарушение приватности \\
unsafe\_instructions & Опасные пошаговые инструкции \\
\bottomrule
\end{longtable}

Уровни риска: \texttt{LOW}, \texttt{MEDIUM}, \texttt{HIGH}.
Решения: \texttt{ALLOW}, \texttt{WARN}, \texttt{BLOCK}.

\section{Требования технического задания}

\begin{enumerate}
    \item REST API на FastAPI со Swagger-документацией.
    \item Два независимых пайплайна модерации: локальный (OSS) и облачный (LLM\,+\,RAG).
    \item Локально: transformer-модель, regex PII, правила jailbreak, risk scoring.
    \item Облачно: ChromaDB, embeddings, RAG, LLM-классификация с JSON-ответом.
    \item Docker, логирование, модульная архитектура, benchmark, pytest, README.
    \item Набор политик (policy dataset) для RAG.
\end{enumerate}

%======================================================================
\chapter{Архитектура решения}
%======================================================================

\section{Общая схема сервиса}

На рисунке~\ref{fig:system} показана высокоуровневая архитектура.

\begin{figure}[H]
\centering
\begin{tikzpicture}[
    node distance=1.1cm and 1.4cm,
    box/.style={rectangle, draw, rounded corners, minimum width=3.2cm, minimum height=0.9cm, align=center, font=\small},
    layer/.style={rectangle, draw, dashed, inner sep=0.35cm},
    arr/.style={-{Stealth[length=2.5mm]}, thick}
]
    \node[box] (client) {Клиент / LLM Gateway};
    \node[box, below=of client] (api) {FastAPI\\GuardrailService};
    \node[box, below left=1.3cm and 0.3cm of api] (local) {LocalModerator\\OSS ensemble};
    \node[box, below right=1.3cm and 0.3cm of api] (cloud) {CloudModerator\\RAG + LLM};
    \node[box, below=of local] (models) {toxic-bert, MiniLM,\\regex, scorer};
    \node[box, below=of cloud] (ext) {ChromaDB,\\OpenRouter};

    \draw[arr] (client) -- (api);
    \draw[arr] (api) -- (local);
    \draw[arr] (api) -- (cloud);
    \draw[arr] (local) -- (models);
    \draw[arr] (cloud) -- (ext);

    \begin{scope}[on background layer]
        \node[layer, fit=(api) (local) (cloud), label={[anchor=north west, font=\footnotesize]x=0.1,y=0.1}Сервис guardrail-service] {};
    \end{scope}
\end{tikzpicture}
\caption{Контекстная архитектура Guardrail Service}
\label{fig:system}
\end{figure}

\section{Слой FastAPI API}

\begin{itemize}
    \item \texttt{app/main.py} --- приложение, CORS, lifespan, маршрутизация \texttt{/api/v1}.
    \item \texttt{app/api/routes.py} --- эндпоинты модерации и health.
    \item \texttt{app/api/schemas.py} --- Pydantic-схемы запроса/ответа.
    \item \texttt{app/services/guardrail\_service.py} --- фасад выбора пайплайна.
\end{itemize}

Валидация входа выполняется на уровне Pydantic (\texttt{content}: 1--32000 символов).
Ошибки бизнес-логики маппятся в HTTP 400/500 с логированием через \texttt{structlog}.

\section{Пайплайн локальной модерации}

Последовательность обработки в \texttt{LocalModerator}:

\begin{enumerate}
    \item Классификация токсичности (\texttt{unitary/toxic-bert}).
    \item Regex-детектор PII и паттернов jailbreak.
    \item Семантический детектор jailbreak (MiniLM, cosine similarity).
    \item Эвристики по ключевым словам (self-harm, malware, fraud и др.).
    \item Агрегация сигналов в \texttt{RiskScorer}.
\end{enumerate}

\section{Пайплайн облачной модерации}

\begin{enumerate}
    \item Формирование запроса: \texttt{content} + \texttt{context}.
    \item RAG: поиск top-$K$ фрагментов политики в ChromaDB.
    \item LLM: классификация с подстановкой политик в промпт.
    \item Нормализация JSON в доменные enum'ы.
\end{enumerate}

\section{RAG и ChromaDB}

Корпус: \texttt{data/policies/policy\_chunks.json} (18 документов).
Эмбеддинги: \texttt{sentence-transformers/all-MiniLM-L6-v2}.
Хранилище: персистентный \texttt{PersistentClient} (\texttt{CHROMA\_PERSIST\_DIR}).
Индексация: \texttt{scripts/ingest\_policies.py}.

\section{Интеграция OpenRouter / OpenAI}

Клиент \texttt{AsyncOpenAI} с параметрами:
\begin{itemize}
    \item \texttt{OPENAI\_API\_KEY} --- ключ OpenRouter или OpenAI;
    \item \texttt{OPENAI\_BASE\_URL} --- по умолчанию OpenAI, для OpenRouter: \texttt{https://openrouter.ai/api/v1};
    \item \texttt{OPENAI\_MODEL} --- идентификатор модели (напр.\ \texttt{openai/gpt-4o-mini}).
\end{itemize}
Ответ запрашивается в формате \texttt{json\_object}, \texttt{temperature=0}.

\section{Логирование и observability}

Используется \texttt{structlog}:
\begin{itemize}
    \item development --- цветной консольный вывод;
    \item production (\texttt{APP\_ENV=production}) --- JSON-логи для агрегации.
\end{itemize}
Ключевые события: \texttt{moderation\_request}, \texttt{local\_moderation\_complete},
\texttt{semantic\_jailbreak\_detected}, \texttt{llm\_classification\_complete},
\texttt{risk\_aggregation\_complete}.

%======================================================================
\chapter{Локальный подход (OSS moderation)}
%======================================================================

\section{Модель toxic-bert}

Модель \texttt{unitary/toxic-bert} (HuggingFace \texttt{pipeline}, text-classification)
возвращает многометочные скоры: \texttt{toxic}, \texttt{severe\_toxic}, \texttt{obscene},
\texttt{threat}, \texttt{insult}, \texttt{identity\_hate}.
Порог \texttt{LOCAL\_MODEL\_THRESHOLD} (по умолчанию 0.65) определяет флаг нарушения.
Маппинг меток на категории риска выполняется в \texttt{transformer.py}.

\textbf{Плюсы:} автономность, низкая стоимость inference после прогрева.
\textbf{Минусы:} слабая чувствительность к jailbreak без явных токсичных маркеров;
ограничение длины входа ($\sim$512 токенов).

\section{Regex PII detector}

\texttt{PIIDetector} ищет шаблоны: email, телефон (US), SSN, номер карты, IPv4,
API-ключи (\texttt{sk-...}), AWS Access Key и др.
При обнаружении формируется сигнал \texttt{pii\_leakage} или \texttt{privacy\_violations}.

\section{Regex jailbreak detection}

Набор \texttt{JAILBREAK\_PATTERNS}: ignore instructions, DAN, role override,
утечка system prompt, delimiter injection, policy bypass и др.
Даёт высокую точность на известных формулировках при низкой латентности.

\section{Семантический детектор jailbreak}

Подробно --- глава~\ref{ch:semantic}. В локальном пайплайне добавляет сигнал
\texttt{semantic\_jailbreak} с категорией \texttt{jailbreak\_attempts}.

\section{Risk scorer и ensemble aggregation}

Каждый детектор порождает \texttt{RiskSignal(source, category, score, detail)}.
Веса категорий заданы в \texttt{CATEGORY\_WEIGHTS} (jailbreak --- 1.0, toxic\_language --- 0.6).
Агрегированный скор:
\[
s_{\mathrm{ens}} = \min\left(\frac{1}{n}\sum_{i=1}^{n} w_i \cdot s_i,\; 1\right)
\]
Уровень риска: $s \geq 0.75 \Rightarrow \mathrm{HIGH}$, $s \geq 0.45 \Rightarrow \mathrm{MEDIUM}$.
Решение: HIGH$\to$BLOCK, MEDIUM$\to$WARN, LOW$\to$ALLOW.

\section{Сводка: плюсы и минусы локального пайплайна}

\begin{table}[H]
\centering
\caption{Локальный пайплайн: оценка}
\begin{tabularx}{\textwidth}{@{}lX@{}}
\toprule
\textbf{Преимущества} & Нет egress в облако LLM; предсказуемая латентность после warm-up; \\
& объяснимые сигналы (source + detail); выше accuracy на тестовом бенчмарке. \\
\midrule
\textbf{Недостатки} & Требует RAM/CPU для моделей; пороги требуют калибровки; \\
& слабее на редких перефразированиях без семантики; нет ссылок на policy ID. \\
\bottomrule
\end{tabularx}
\end{table}

%======================================================================
\chapter{Облачный подход (LLM + RAG)}
%======================================================================

\section{Retrieval pipeline}

\begin{lstlisting}[style=python]
# PolicyRAG.retrieve(query) -> list[dict]
results = collection.query(query_texts=[query], n_results=top_k)
\end{lstlisting}

Запрос --- конкатенация \texttt{context} и \texttt{content}.
Метаданные чанков: \texttt{category}, \texttt{risk\_level}, \texttt{title}.

\section{Embeddings}

Единая модель \texttt{all-MiniLM-L6-v2} для ChromaDB и семантического jailbreak.
Размерность 384, cosine distance в HNSW-индексе Chroma.

\section{Policy chunks}

18 JSON-объектов с полями \texttt{id}, \texttt{title}, \texttt{category},
\texttt{risk\_level}, \texttt{text}. Покрывают все категории ТЗ.

\section{Structured JSON classification}

Системный промпт фиксирует схему ответа LLM. Пример полей:
\texttt{decision}, \texttt{category}, \texttt{risk\_level}, \texttt{explanation}.
Парсинг через \texttt{json.loads}, нормализация в \texttt{\_normalize()}.

\section{Retry logic}

Декоратор \texttt{tenacity}: до 3 попыток, экспоненциальная задержка 1--8\,с.
Снижает влияние сетевых сбоев OpenRouter.

\section{Fallback strategy}

При ошибке парсинга JSON: \texttt{decision=WARN}, \texttt{category=none},
\texttt{risk\_level=MEDIUM} --- консервативная стратегия (не пропускать сомнительный контент молча).

\section{Плюсы и минусы облачного пайплайна}

\begin{table}[H]
\centering
\caption{Облачный пайплайн: оценка}
\begin{tabularx}{\textwidth}{@{}lX@{}}
\toprule
\textbf{Преимущества} & Гибкость формулировок политики; естественные объяснения; \\
& быстрое обновление политик через re-ingest без деплоя правил. \\
\midrule
\textbf{Недостатки} & Зависимость от API; выше латентность и стоимость; \\
& риски трансграничной передачи данных; на бенчмарке accuracy 5/10. \\
\bottomrule
\end{tabularx}
\end{table}

%======================================================================
\chapter{Semantic Jailbreak Detection}
\label{ch:semantic}
%======================================================================

\section{Назначение}

Regex ловит точные и близкие шаблоны, но пропускает перефразирования
(``обойди ограничения безопасности другими словами'').
Семантический модуль сопоставляет вход с \emph{каноническими} jailbreak-примерами
в пространстве эмбеддингов.

\section{sentence-transformers}

Модель \texttt{SemanticJailbreakDetector} загружает \texttt{SentenceTransformer}
один раз на процесс (\texttt{@lru\_cache}). Кодирование выполняется под \texttt{torch.no\_grad()}.

\section{Cosine similarity}

Векторы L2-нормализуются (\texttt{normalize\_embeddings=True}).
Косинусное сходство реализовано скалярным произведением:
\[
\cos(\mathbf{q}, \mathbf{c}_i) = \mathbf{q}^\top \mathbf{c}_i
\]
Берётся $\max_i \cos(\mathbf{q}, \mathbf{c}_i)$; при $\max \geq \tau$ --- детекция.

\section{Canonical jailbreak prompts}

23 эталонных фразы (DAN, ignore instructions, grandma exploit, STAN mode и др.)
хранятся в \texttt{CANONICAL\_JAILBREAKS}. Эмбеддинги кэшируются при первом вызове
\texttt{\_encode\_canonical()}.

\section{Threshold tuning}

Параметр \texttt{SEMANTIC\_JAILBREAK\_THRESHOLD} (по умолчанию 0.62).
Повышение $\tau$ снижает false positive, но увеличивает false negative.
Рекомендуется калибровка на валидационной выборке заказчика.

\section{Отличие от regex}

\begin{table}[H]
\centering
\caption{Regex vs семантический детектор}
\begin{tabular}{@{}lll@{}}
\toprule
 & \textbf{Regex} & \textbf{Semantic} \\
\midrule
Механизм & Подстрока / паттерн & Сходство в embedding space \\
Перефразы & Слабо & Сильнее \\
Ложные срабатывания & На случайные совпадения & На близкие по смыслу безопасные фразы \\
Латентность & Микросекунды & Десятки--сотни мс (encode) \\
\bottomrule
\end{tabular}
\end{table}

Интерфейс: \texttt{detect(text) -> (bool, float, str|None)} --- флаг, скор, matched example.

%======================================================================
\chapter{API и интерфейсы}
%======================================================================

\section{Эндпоинты}

\begin{table}[H]
\centering
\caption{REST API}
\begin{tabular}{@{}llp{6.5cm}@{}}
\toprule
\textbf{Метод} & \textbf{Путь} & \textbf{Назначение} \\
\midrule
GET & /api/v1/health & Статус и версия \\
POST & /api/v1/moderate & Модерация (approach в теле) \\
POST & /api/v1/moderate/local & Принудительно local \\
POST & /api/v1/moderate/cloud & Принудительно cloud \\
GET & /docs & Swagger UI \\
\bottomrule
\end{tabular}
\end{table}

\section{Схема запроса и ответа}

\begin{lstlisting}[style=json]
{
  "content": "Текст ответа ассистента",
  "approach": "local",
  "context": "Опциональный контекст диалога"
}
\end{lstlisting}

\begin{lstlisting}[style=json]
{
  "decision": "BLOCK",
  "category": "jailbreak_attempts",
  "risk_level": "HIGH",
  "explanation": "Ensemble moderation detected...",
  "approach": "local",
  "confidence": 0.87,
  "retrieved_policies": null
}
\end{lstlisting}

\section{Healthcheck}

\texttt{GET /api/v1/health} возвращает \texttt{status}, \texttt{version}, список \texttt{approaches}.
Используется в Docker HEALTHCHECK.

%======================================================================
\chapter{Benchmark и сравнение подходов}
%======================================================================

Датасет: \texttt{tests/benchmark\_dataset.json}, 10 размеченных примеров.
Скрипт: \texttt{scripts/benchmark\_compare.py}.

\section{Результаты}

\begin{table}[H]
\centering
\caption{Сравнение пайплайнов на тестовом наборе}
\begin{tabular}{@{}lcc@{}}
\toprule
\textbf{Метрика} & \textbf{Local (OSS)} & \textbf{Cloud (LLM+RAG)} \\
\midrule
Accuracy & 7 / 10 (70\%) & 5 / 10 (50\%) \\
Средняя латентность & 1402 ms & 3515 ms \\
\bottomrule
\end{tabular}
\end{table}

\section{Trade-off анализ}

\begin{longtable}{@{}p{3.2cm}p{5cm}p{5cm}@{}}
\toprule
\textbf{Критерий} & \textbf{Local} & \textbf{Cloud} \\
\midrule
\endhead
Латентность & Ниже (после warm-up) & Выше (сеть + LLM) \\
Стоимость & Только compute & Токены + API \\
Explainability & Источник сигнала (transformer, regex, semantic) & Текст LLM + policy IDs \\
Jailbreak robustness & Regex + semantic + rules & Зависит от промпта и RAG \\
Масштабируемость & Горизонтально, stateless* & Rate limits провайдера \\
\bottomrule
\multicolumn{3}{@{}p{\linewidth}@{}}{\footnotesize *Модели в памяти процесса; для 2M msg/h нужен пул воркеров.}
\end{longtable}

%======================================================================
\chapter{Production-рекомендация}
%======================================================================

\section{Local как основной контур}

На основании бенчмарка и архитектурных свойств \textbf{рекомендуется}
использовать локальный пайплайн для inline-модерации каждого ответа:
\begin{itemize}
    \item выше accuracy на фиксированном наборе;
    \item ниже латентность;
    \item отсутствие передачи текста во внешний LLM.
\end{itemize}

\section{Cloud как secondary layer}

Облачный пайплайн целесообразен для:
\begin{itemize}
    \item случаев \texttt{WARN} (второе мнение);
    \item асинхронной очереди human/LLM review;
    \item сценариев, где критична ссылка на текст корпоративной политики.
\end{itemize}

\section{Гибридная архитектура}

\begin{figure}[H]
\centering
\begin{tikzpicture}[
    box/.style={rectangle, draw, rounded corners, minimum width=2.8cm, minimum height=0.8cm, align=center, font=\small},
    arr/.style={-{Stealth}, thick}
]
    \node[box] (out) {Ответ LLM};
    \node[box, right=1.2cm of out] (loc) {Local};
    \node[box, right=1.2cm of loc] (dec) {Решение};
    \node[box, below=1cm of dec] (cloud) {Cloud (async)};

    \draw[arr] (out) -- (loc);
    \draw[arr] (loc) -- node[above, font=\footnotesize]{ALLOW/BLOCK} (dec);
    \draw[arr] (loc) -- node[left, font=\footnotesize]{WARN} (cloud);
    \draw[arr] (cloud) -- (dec);
\end{tikzpicture}
\caption{Рекомендуемый гибридный контур}
\end{figure}

%======================================================================
\chapter{Соответствие требованиям 152-ФЗ}
%======================================================================

Федеральный закон №\,152-ФЗ регулирует обработку персональных данных (ПДн)
граждан РФ. Модерация может обрабатывать ПДн из пользовательских сообщений.

\section{Риски облачного LLM}

Передача текста в OpenRouter/зарубежного провайдера может constitute
трансграничную передачу ПДн. Требуется правовое основание, уведомление РКН,
договор с оператором, оценка страны получателя.

\section{Data residency}

Для контуров с требованием локализации --- размещение сервиса и логов в РФ,
\textbf{отключение cloud moderation} для потоков с ПДн, использование только local pipeline.

\section{Логирование и PII}

\texttt{structlog} не должен писать полный пользовательский текст на INFO в prod.
Рекомендации: маскирование, хеширование, сокращённый preview, отдельный audit storage
с контролем доступа и сроком хранения.

%======================================================================
\chapter{Масштабирование и production concerns}
%======================================================================

\section{Целевая нагрузка: 2 млн сообщений/час}

$\approx 556$ RPS среднее, пики $\times 3$--$5$. Один синхронный запрос local
$\sim$50--200\,ms после warm-up (бенчмарк 1402\,ms включает cold start моделей).

\section{Стратегии}

\begin{itemize}
    \item \textbf{Horizontal scaling} --- N реплик FastAPI за балансировщиком.
    \item \textbf{Model warmup} --- прогрев toxic-bert и MiniLM при старте пода.
    \item \textbf{Caching} --- кэш эмбеддингов канонических jailbreak; опционально LRU по hash(content).
    \item \textbf{Async processing} --- WARN/BLOCK review через очередь (Kafka/RabbitMQ).
    \item \textbf{Rate limiting} --- на API gateway; защита от abuse.
    \item \textbf{Observability} --- метрики Prometheus: \texttt{moderation\_latency\_seconds},
    \texttt{decision\_total\{decision,category\}}, \texttt{semantic\_hits\_total};
    дашборды Grafana, алерты на рост BLOCK/WARN.
\end{itemize}

%======================================================================
\chapter{Тестирование}
%======================================================================

\begin{table}[H]
\centering
\caption{Покрытие тестами}
\begin{tabular}{@{}lp{8.5cm}@{}}
\toprule
\textbf{Модуль} & \textbf{Содержание} \\
\midrule
\texttt{test\_api.py} & Health, local moderate, PII, validation 422 \\
\texttt{test\_jailbreak.py} & 10+ jailbreak/PII/unsafe кейсов \\
\texttt{test\_semantic\_detector.py} & Кэш, порог, safe vs jailbreak \\
\texttt{benchmark\_dataset.json} & 10 labeled samples для benchmark \\
\bottomrule
\end{tabular}
\end{table}

Запуск: \texttt{pytest -v}. Cloud-тесты требуют API-ключа и опциональны в CI.

%======================================================================
\chapter{Заключение}
%======================================================================

\section{Реализовано}

\begin{itemize}
    \item Production-style MVP: FastAPI, два пайплайна, Docker, Swagger.
    \item Локальный ensemble: toxic-bert, PII, regex jailbreak, semantic jailbreak, scorer.
    \item Облачный контур: ChromaDB RAG, OpenRouter-совместимый LLM, retry, JSON fallback.
    \item 18 категорий, structured output, benchmark, pytest.
\end{itemize}

\section{Направления улучшения}

Калибровка порогов на расширенном датасете; гибридный score local+cloud;
ONNX/GPU; версионирование политик; Prometheus middleware; мультиязычные эталоны jailbreak.

\section{Итоговая рекомендация}

Для production-контура conversational AI с упором на безопасность, латентность
и соответствие 152-ФЗ: \textbf{основной путь --- LocalModerator};
\textbf{CloudModerator} --- вторичный слой для спорных случаев и policy-grounded review.

\end{document}
